\documentstyle[% 


righttag% 


,11pt% 


,amssymb% 


,russian%               remove in Eng. 


,izvru%                 Set izven% in Eng. 


]{amsart} 


 


%  {\in}


\newcommand{\tts}[1]{} 


 


%\hoffset=-15mm%                  For Eng. 


\setcounter{page}{43} 


\god{2001} 


\nomer{7~(470)} %                        Remove in Eng. 


%\nomer{Vol.\,45, No.\,5}%               Set in Eng. 


%\str{\pageref{begin}--\pageref{end}}%  Set in Eng. 


\udk{517.98} 


 


\author{\it Ш.М.\,Усманов} 


% NB:   ^^ remove in Eng. 


\title{Условные ожидания на вещественных  


$W^*$-алгебрах\\ и $JW$-алгебрах} 


 


\date{} 


%\pagestyle{myheadings}%         Set in Eng. 


\pagestyle{plain} %              Remove in Eng. 


\begin{document} 


%\thispagestyle{plain}%          Set in Eng. 


\maketitle 


\label{begin} 


 


Условным ожиданием на $W^*$-алгебре называется линейное положительное 


нормальное отображение, являющееся проекцией (т.\,е. идемпотентное) 


и имеющее норму 1. В [1] Такесаки дал необходимые 


и достаточные условия существования условного ожидания на 


$W^*$-алгебре, сохраняющего нормальный полуконечный вес. 


В [2]--[5] были 


введены условные ожидания на йордановых операторных алгебрах 


($JW$-алгебрах) и исследовалась их связь с условными ожиданиями на 


их обертывающих $W^*$-алгебрах. 


В [6] для $JW$-алгебр был доказан один вариант теоремы 


Такесаки о существовании условных ожиданий, сохраняющих


точный нормальный полуконечный вес. Затем в [7] для 


$JBW$-алгебр был получен аналог теоремы Такесаки о существовании 


условных ожиданий, сохраняющих точный нормальный линейный 


функционал единичной нормы (состояние). 


 


Цель данной работы --- получение необходимого и достаточного 


условия существования условного ожидания на обратимой 


$JW$-алгебре, сохраняющего точный нормальный полуконечный вес. 


\medskip 


 


{\bf 1.} {\it Предварительные сведения.} $JW$-алгеброй называется 


вещественное 


линейное пространство $A$ самосопряженных операторов из $B(H)$, 


замкнутое относительно йорданова (симметризованного) умножения 


$x\circ y=\frac{1}{2}(xy+yx)$, $x,y\in A$, и замкнутое в слабой 


операторной топологии.\; 


$JW$-алгебра $A$ называется обратимой, если $x_1x_2\dotsc x_n+x_n 


x_{n-1}\dotsc x_1{\in} A$ 


для всех $x_1,x_2,\dotsc,x_n{\in} A$; такие $JW$-алгебры характеризуются 


тем, 


что $A=R(A)_{sa}=\{x\in R(A):x=x^*\}$, где $R(A)$ --- наименьшая слабо 


замкнутая вещественная подалгебра $B(H)$, содержащая $A$. Обратимая 


$JW$-алгебра $A$ разлагается по центру в прямую сумму $JW$-алгебр 


$A_r$ и $A_c$: $A=A_r\oplus A_c$, где $A_c$ --- множество 


самосопряженных операторов 


$W^*$-алгебры $R(A_c)=\cal U(A_c)$, и $A_r$ --- множество 


самосопряженных 


операторов вещественной $W^*$-алгебры $R(A_r)$. 


 


Напомним, что вещественной $W^*$-алгеброй называется вещественная 


$*$-алгебра $R$ ограниченных линейных операторов в комплексном 


гильбертовом пространстве $H$, содержащая единицу, замкнутая в слабой 


операторной топологии и такая, что $R\cap iR=\{0\}$. Обертывающая 


$W^*$-алгебра $\cal U(R)$ для $R$ (т.\,е. наименьшая $W^*$-алгебра, 


содержащая $R$) 


имеет вид: $\cal U(R)=R+iR$, и $R$ определяет на $\cal U(R)$ некоторый 


инволютивный 


$*$-антиавтоморфизм $\alpha_R:\alpha_R(x+iy)=x^*+iy^*$, где $x,y\in R$. 


Соответственно $R$ определяется при помощи $\alpha_R:R=\{x\in\cal 


U(R):\alpha_R(x)=x^*\}$. 


 


О теории $JW$-алгебр см. [8]--[10], о теории вещественных 


$W^*$-алгебр см. [11]. 


\medskip 


 


{\bf 2.} {\it Условные ожидания на вещественных $W^*$-алгебрах.} 


Определим 


условное ожидание на вещественной $W^*$-алгебре относительно ее 


вещественной $W^*$-подалгебры. 


\medskip 


 


{\bf Определение.} Пусть $R$ --- вещественная $W^*$-алгебра, $Q$ --- ее 


вещественная $W^*$-подалгебра. Линейное отображение $E:R\to Q$ 


называется 


условным ожиданием, если $E$ является положительной проекцией 


с нормой 1.\; $E$ называется нормальным, если $x_i\nearrow x$ влечет 


$T(x_i)\nearrow T(x)$ для $\{x_i\}_{i\in Z}$, $x$ из $R_s$; $E$ 


называется точным, если 


$T(x^*x)=0$ влечет $x=0$. 


\medskip 


 


Сформулируем теорему о существовании условного ожидания на 


вещественной $W^*$-алгебре. 


 


\begin{thm} 


Пусть $R$, $Q$ --- вещественные $W^*$-алгебры, $Q\subset R$, $\varphi$ --- 


точный нормальный полуконечный вес на $R^+$. Следующие условия 


эквивалентны$:$ 


\begin{itemize} 


\item[1)] существует точное нормальное условное ожидание $E:R\to Q$ 


такое, что 


$$ 


\varphi(x)=\varphi(E(x))\ \; \forall x\in R; 


$$ 


\item[2)] точный нормальный вес $\varphi|_{Q^+}$ является полуконечным 


и 


$$ 


\sigma^{\ov\varphi}_t(\cal U(Q))=\cal U(Q)\ \; \forall t\in\bold R 


$$ 


для веса $\ov\varphi$, являющегося $\alpha_Q$-инвариантным продолжением 


$\varphi$ на $\cal U(R^+)$. 


\end{itemize} 


\end{thm} 


 


\begin{pf} 


$(1)\ \Rightarrow\ (2)$. Пусть $E:R\to Q$ --- $\varphi$-инвариантное 


условное ожидание, $E_0$ --- его сужение на $JW$-алгебру $R_s$; тогда из 


положительности $E$ следует, что $E_0(R_s)\subset Q_s$. По основной 


теореме 


и следствию 1 работы [3] существует точное нормальное условное 


ожидание $\ov E:\cal U(R_s)\to\cal U(Q_s)$ такое, что $E_0=E_0\cdot\ov 


E$. Пусть $\ov\varphi$ --- 


вес на $M^+$, являющийся $\alpha_R$-инвариантным продолжением 


$\varphi$, $\alpha_r$ --- 


инволютивный $*$-антиавтоморфизм $M$, определяющий $R$. Тогда по 


построению $\ov E$ переводит самосопряженные элементы в 


самосопряженные, 


а кососопряженные элементы в кососопряженные, и тогда для $\ov E$


выполняется $\ov\varphi\ov E=\ov\varphi$. Действительно, 


$$ 


\ov\varphi(\ov E(x+iy))=\ov\varphi(\ov E(x)+ 


\ov E(iy))=\varphi(E(x))+i\varphi(E(y))= 


\varphi(x)=\ov\varphi(x+iy) 


$$ 


для всех $x+iy\in M^+$, $x\in R^+$, $y^+=-y$, т.\,к. в силу 


$\alpha_R$-инвариантности 


для веса $\varphi$ выполняется $\varphi(y)=\varphi(\ov E(y))=0$. 


Тогда по теореме Такесаки $\sigma^{\ov\varphi}_t(N)=N$ для всех 


$t\in\bold R$. 


\smallskip 


 


$(2)\ \Rightarrow\ (1)$. В силу условия 2) по теореме Такесаки 


существует точное 


нормальное условное ожидание $D:M\to N$ такое, что 


$\ov\varphi{}D=\ov\varphi$. 


Покажем, что $D$ коммутирует с $\alpha_R$. 


 


Построим представление Гельфанда--Наймарка--Сигала 


$\pi_{\ov\varphi}(M)$\; 


$W^*$-алгебры $M$ по весу $\ov\varphi$; тогда гильбертово пространство 


$H_\varphi$ является 


пополнением $\{x\in M:\ov\varphi(x^*x)<+\infty\}$ по норме 


$\ov\varphi(x^*x)^{1/2}$, и, кроме 


того, имеет вид $H=K+i{}K$, где $K$ --- вещественное гильбертово 


пространство, являющееся пополнением по той же норме множества 


$\Re_{\ov\varphi}=\{x\in R:\ov\varphi(x^*x)<+\infty\}$. Легко 


проверяется, что в силу $\alpha_R$-инвариантности $\ov\varphi$ 


соответствующее скалярное произведение $(\ ,\ )_{\ov\varphi}$ 


принимает вещественные значения на $K$ и из $R\cap iR=\{0\}$ следует 


$K\cap i{}K=\{0\}$. 


 


Пусть $H_1$ --- гильбертово подпространство $H_{\ov\varphi}$, 


порожденное 


множеством $\Re_{\ov\varphi|_N}=\{x\in N:\ov\varphi(x^*x)<+\infty\}$, 


$p:H_{\ov\varphi}\to H_1$ --- 


ортогональный 


проектор, порождающий условное ожидание~$D$~[1], 


$$ 


D(x)=\pi^{-1}_{\ov\varphi}(p\pi_{\ov\varphi}(x)p),\quad x\in M. 


$$ 


Кроме того, рассмотрим вещественное подпространство 


$K_1\subset K$, порожденное множеством $\Re_{\varphi|_Q}=\{x\in 


Q:\ov\varphi(x^*x)<+\infty\}$; 


очевидно, $H_1=K+i{}K$. Пусть $p_1$ --- ортогональный проектор из $K$ 


на 


$K_1$. Так как $K\cap i{}K=\{0\}$, то $p_1$ можно продолжить по 


линейности 


до ортогонального проектора $p'$ из $H_{\ov\varphi}$ на 


$H_1=K_1+i{}K_1$, и в силу 


единственности ортогонального проектора на подпространство $H_1$ 


проектор $p'$ совпадает с $p$. 


 


Таким образом, условное ожидание $D$ коммутирует с $\alpha_R$ и 


сужается 


до условного ожидания $E:R\to Q$,


$$ 


E(x)=\pi_{\ov\varphi}(p_1\pi_{\ov\varphi}(x)p_1),\quad x\in R. 


$$ 


Единственность $E$ следует из теоремы Такесаки, выполнение условия 


$$ 


\varphi(x)=\varphi(E(x))\ \; \forall x\in R 


$$ 


очевидно по построению $E$. 


\end{pf} 


 


{\bf 3.} {\it Условные ожидания на обратимых $JW$-алгебрах.} Условным 


ожиданием 


на $JW$-алгебре $A$ относительно ее $JW$-подалгебры $A_1$ называется 


нормальное линейное положительное идемпотентное отображение $\Phi:A\to 


A_1$, 


имеющее норму~1 [3]. 


 


Напомним, что в [9] была построена теория Томиты--Такесаки 


для состояний на $JBW$-алгебрах; развивая ее, Н.\,Трунов [12], [13] 


построил 


аналогичную теорию для весов на $JBW$-алгебрах. В 


рамках этой теории для нормального состояния или нормального 


полуконечного веса $\varphi$ на $JBW$-алгебре $A$ был определен аналог 


группы 


модулярных автоморфизмов --- косинус-группа 


$\{\rho^\varphi_t\}_{t\in\bold R}$ слабо непрерывных 


положительных линейных отображений на $A$, сохраняющих единичный 


оператор и таких, что $\varphi(\rho^\varphi_t(x)\circ y)=\varphi(x\circ 


\rho^\varphi_t(y))$ 


для всех $x,y\in A$, $t\in\bold R$, и выражение 


$$ 


s(x,y)=\int^{+\infty}_{-\infty} 


\varphi(x\circ \rho^\varphi_t(y))\text{\,cosh\,}(\pi t)^{-1}dt 


$$ 


определяет самополярную форму, ассоциированную с $\varphi$. Для 


обратимой $JW$-алгебры $A$ таким косинусным семейством является 


семейство 


$$ 


\rho^\varphi_t=\frac{1}{2}(\sigma^{\ov\varphi}_t+ 


\sigma^{\ov\varphi}_{-t})\big|_A, 


$$ 


где $\ov\varphi$ --- вес, являющийся $\alpha_A$-инвариантным 


продолжением $\varphi$ на $\cal U(A)^+$, 


именно такой, что $\ov\varphi(a+i{}b)=\varphi(a)$ для $a+i{}b\in\cal 


U(A)^+$, $a\ge0$, $b^*=-b$; 


$\alpha_A$ --- инволютивный $*$-антиавтоморфизм $\cal U(A)$, 


порождающий $R(A)$ в $\cal U(A)$. 


 


Хагеруп и Штермер ([7], теорема 4.2) доказали, что 


для $JW$-алгебры $A$ с точным 


нормальным состоянием $\varphi$ и ее $JW$-подалгебры $A_1$ с состоянием 


$\psi=\varphi|_{A_1}$ 


существует точное нормальное условное ожидание $E:M\to N$ такое, 


что $\varphi=\psi\cdot E$ тогда и только тогда, когда 


$$ 


\rho^\psi_t(x)=\rho^\varphi_t(x),\quad 


x\in A_1,\ \; t\in\bold R. 


$$ 


 


Следующее предложение обобщает следствие 4.3 из [7]. 


\medskip 


 


{\bf Предложение.} {\it Пусть $M$ --- $W^*$-алгебра, $A$ --- обратимая 


$JW$-подалгебра 


$M$ такая, что $Z(A)=Z(\cal U(A)_{sa})$, $\varphi$ --- точный 


нормальный 


полуконечный $\alpha$-инвариантный вес на $M^+$, для которого 


$$ 


\rho^\varphi_t(x)\in A\ \; \forall t\in\bold R, 


\ \; \forall x\in A. 


$$ 


Тогда $\sigma^\varphi_t(\cal U(A))=\cal U(A)$, $t\in\bold R$. 


} 


 


\begin{pf} 


Для случая, когда $\varphi$ --- ограниченный линейный 


функционал, утверждение предложения доказано в следствии 4.3 [7]. 


Сведем наше доказательство к этому случаю. 


 


Предположим, что существуют $e\in A$, $t_0\in\bold R$ такие, что 


$\sigma^\varphi_{t_0}(e)\notin\cal U(A)$. 


Так как по спектральной теореме для $JW$-алгебр (см. [9], [10] или 


[11]) всякий элемент $JW$-алгебры можно разложить в слабо сходящийся 


интеграл по спектральному семейству проекторов, то можем 


без ограничения общности взять проектор $e\in A$, а в силу 


полуконечности 


$\varphi$ можем выбрать $e$ таким, что $\varphi(e)<+\infty$. Пусть 


$\sigma^\varphi_{t_0}(e)=p+i{}q$, $q\ne0$. Пусть $g=s_{\cal 


U(A)_\varphi}(e)$ --- относительный носитель $e$ 


в $W^*$-алгебре $\cal U(A)_\varphi$, которая является централизатором 


веса $\varphi$, 


$$ 


\cal U(A)_\varphi=\{x\in\cal U(A):\sigma^\varphi_t(x)= 


x\ \forall t\in\bold R\}. 


$$ 


Как и всякий центральный носитель в $W^*$-алгебре, $g$ имеет вид 


$g=\cup\{e'\in\cal U(A)_\varphi:e'\preceq e\}$. 


Так как $e\in A$, то $a_A(e)=e$, откуда если $e'\preceq e$, то 


$\alpha_A(e')\preceq e$, и тогда $g\in\cal U(A)_\varphi\cap A$. Так как 


вес $\varphi$ является точным нормальным 


$\alpha_A$-инвариантным полуконечным следом на $\sigma$-конечной 


$W^*$-алгебре $\cal U(A)_\varphi$, то существует последовательность 


$\{g_n\}^\infty_{n=1}$ попарно 


ортогональных проекторов из $A\cap\cal U(A)_\varphi$ таких, что 


$g=\sum\limits^\infty_{n=1}g_n$, и, 


кроме того, $\varphi(g_n)<\infty$ для всех $n$. Разложим проектор $g$ 


по $\{g_n\}$ 


$$ 


e=\sum^\infty_{n=1}e_{n,m}, 


$$ 


где $e_{n,n}=g_n e g_n$, $e_{n,m}=g_n e g_m+g_m e g_n$ при $n\ne m$. 


Так как проекторы $g_n$ 


выбраны из $\cal U(A)_\varphi$, то 


$$ 


g_n(p+i q)g_m=g_n(\sigma^\varphi_{t_0}(e))g_m= 


\sigma^\varphi_{t_0}(g_n e g_m). 


$$ 


Так как $p+i{}q\notin\cal U(A)$, то существуют числа $n_0$, $m_0$ 


такие, что $g_{n_0}(p+i{}q)g_{m_0}=\sigma^\varphi_{t_0}(g_{n_0}e 


g_{m_0})\notin\cal U(A)$. 


Пусть теперь $r=g_{n_0}+g_{m_0}$, $s=rer$; тогда 


для точного нормального ограниченного линейного функционала 


$\psi=\varphi\big|_{\cal U(A)_\varphi}$ 


на $W^*$-алгебре $r(\cal U(A))r$ выполняется 


$$ 


\sigma^\psi_t(x)+\sigma^\psi_{-t}(x)= 


\sigma^\varphi_t(r x r)+\sigma^\varphi_{-t}(r x r)= 


r(\sigma^\varphi_t(x))r+r(\sigma^\varphi_{-t}(x))r 


\in r A r 


$$ 


для всех $t\in\bold R$, $x\in r A r$; одновременно 


$\sigma^\psi_{t_0}(rer)=r(\sigma^\varphi_{t_0}(e))r\notin 


r(\cal U(A))r$. 


Получили противоречие с утверждением следствия 4.3 из [7]. 


\end{pf} 


 


Теперь можем сформулировать необходимое и достаточное условие 


существования условного ожидания на обратимой $JW$-алгебре в 


терминах модулярных косинус-групп. 


 


\begin{thm} 


Пусть $A$ --- обратимaя $JW$-алгебра, $A_1$ --- ее $JW$-подалгебра. 


Следующие условия эквивалентны. 


\begin{itemize} 


\item[1)] На $A^+$ существует точный нормальный полуконечный вес 


$\varphi$ 


такой, что его сужение $\psi$ на $A^+_1$ является точным нормальным 


полуконечным 


весом и для косинус-группы $\{\rho^\varphi_t\}_{t\in\bold R}$ 


выполняется 


$\rho^\varphi_t(A_1)=A_1$ 


для всех $t\in\bold R$. 


\item[2)] Существует тoчное нормальное условное ожидание $T:A\to A_1$ 


такое, что $\psi{}T(x)=\psi(x)$ для всех $x\in A$. 


\end{itemize} 


\end{thm} 





\begin{pf} 


$(1)\ \Rightarrow\ (2)$. Следует из предложения 3 и теоремы 1. 


\smallskip 


 


$(2)\ \Rightarrow\ (1)$. В силу основной теоремы и следствия 1 из [3] 


существует 


точное нормальное ожидание $D:\cal U(A)\to A''_1=\cal U(A_1)$, 


продолжающее 


$T$ и такое, что $D\alpha_A=\alpha_A{}D$. Рассмотрим вес~$\varphi$, 


являющийся 


$\alpha_A$-инвариантным продолжением $\psi$ на $\cal U(A)^+$; тогда 


$$ 


\psi=\frac{\varphi+\varphi\alpha_A}{2}


$$ 


и для любого $x\in\cal U(A)$ выполняется 


\begin{multline*} 


\varphi(D(x))=\psi\bigg( 


\frac{D(x)+\alpha_A(D(x))}{2}\bigg)=\psi\bigg( 


\frac{D(x)+D(\alpha_A(x))}{2}\bigg)= \\ 


% 


=\psi\bigg(D\bigg(\frac{x+\alpha_A(x)}2{} \bigg) \bigg)= 


\psi\bigg(T\bigg(\frac{x+\alpha_A(x)}{2} \bigg) \bigg)= 


\psi\bigg(\frac{x+\alpha_A(x)}{2} \bigg)=\varphi(x). 


\end{multline*} 


Итак, всякое условное ожидание $T:A\to A_1$, удовлетворяющее условиям 2) 


теоремы, можно продолжить до условного ожидания $D:\cal U(A)\to\cal 


U(A_1)$ такого, что 


$\varphi(D(x))=\varphi(x)$ $\forall x\in\cal U(A)$. 


Тогда по теореме Такесаки имеем 


$\sigma^\varphi_t(\cal U(A_1))=\cal U(A_1)$ 


$\forall t\in\bold R$, 


откуда следует утверждение~1) теоремы. 


\end{pf} 


 


 


\begin{thebibliography}{99} 


 


\bibitem{1} 


Takesaki M. {\em Conditional expectation in von Neumann algebra} 


// J. Funct. Anal. -- 1972. -- V.\,9. -- \No\,3. -- P.\,306--321. 


\bibitem{2} 


Edwards C.M. {\em Conditional expectations on Jordan algebras} 


// Fund. Aspects of 


Quantum Theory, NATO ASI Ser. -- Plenum Press, 1985. -- P.\,75. 


\bibitem{3} 


Stormer E. {\em Positive projections onto Jordan algebras and 


their enveloping von Neumann algebras} // Ideas and Methods 


in Math. Anal., Stochastics and Appl. -- Cambridge 


Univ. Press, 1992. -- P.\,389--393. 


\bibitem{4} 


Аюпов Ш.А. {\em Условные математические ожидания и мартингалы на 


йордановых алгебрах} // ДАН УзССР. -- 1981. -- \No\,10. -- С.\,3--5. 


\bibitem{5} 


Аюпов Ш.А. {\em Классификация, представления и вероятностные аспекты 


упорядоченных йордановых алгебр}. Дис. $\dotsc$ докт. физ.-матем. наук. 


-- 


Ташкент, 1982. -- 295 с. 


\bibitem{6} 


Аюпов Ш.А., Азизов Э., Усманов Ш.М. {\em О существовании условных 


ожиданий в йордановых операторных алгебрах} // Узбекск. матем. журн. 


-- 1991. -- \No\,6. -- С.\,11--14. 


\bibitem{7} 


Haagerup U., Stormer E. {\em Positive projections of von Neumann 


algebras onto $JW$-algebras} // J. Reports Math. Phys. -- 1995. 


-- V.\,36. -- P.\,317--330. 


\bibitem{8} 


Аюпов Ш.А. {\em Классификация и представление упорядоченных йордановых 


алгебр}. -- Ташкент: Фан, 1986. -- 123 с. 


\bibitem{9} 


Hanche-Olsen H., Stormer E. {\em Jordan operator algebras}. -- Boston 


e.\,a.: Pitman, 1984. -- 183 p. 


\bibitem{10} 


Topping D. {\em Jordan algebras of self-adjoint operators} 


// Memoirs Amer. Math. Soc. -- 1965. -- \No\,53. -- P.\,1--48. 


\bibitem{11} 


Ayupov Sh., Rakhimov A., Usmanov Sh. {\em Jordan, real and Lie 


structures in operator algebra} // Ser.\,MIA \No\, 418. 


-- Dordrecht: Kluwer AP, 1998. -- 235 р. 


\bibitem{12} 


Трунов Н.В. {\em Интегрирование относительно веса на йордановых 


алгебрах}. -- Казанский ун-т. -- Казань, 1984. -- 34 с. 


Деп. в ВИНИТИ 10.07.1984, \No\,4948--84. 


\bibitem{13} 


Трунов Н.В., Шерстнев А.Н. {\em Введение в теорию некоммутативного 


интегрирования} // Итоги науки и техн. ВИНИТИ. Современ. пробл. матем. 


-- 1985. -- Т.\,27. -- С.\,167--190. 


 


\end{thebibliography} 


 


\vskip 2cm 


 


\post{\it Институт математики Академии}{\it Поступила} 


\post{\it наук Республики Узбекистан}{19.01.1999} 


 


\label{end} 


\end{document} 


 


Haagerup U., Hanche-Olsen H. {\em A Tomita--Takesaki theory for 


Jordan algebras} // 


J. Oper. Theory. -- 1984. -- V.\,11. -- \No\,2. -- P.\,343--364. 


 


 


